\documentclass[a4paper,12pt]{article}

\usepackage[T2A]{fontenc}
\usepackage[utf8]{inputenc}
\usepackage[russian]{babel}

\usepackage{amsmath, amssymb, amsthm, mathtools}
\usepackage{helvet}
\renewcommand{\familydefault}{\sfdefault}
\usepackage{icomma}
\usepackage{geometry}
\usepackage{microtype}
\usepackage{tikz}
\usepackage{tkz-euclide}
\usepackage{pgfplots}
\pgfplotsset{compat=1.18}
\usepackage{tabularx, booktabs, longtable}
\usepackage{enumitem}
\usepackage{hyperref}
\usepackage{parskip}
\usepackage{fancyhdr}
\usepackage{setspace}
\usepackage{xcolor}

\geometry{left=18mm,right=18mm,top=18mm,bottom=20mm}
\hypersetup{hidelinks}
\onehalfspacing

\definecolor{accent}{HTML}{008080}
\definecolor{darkaccent}{HTML}{004D4D}
\definecolor{textgray}{HTML}{333333}
\definecolor{softgray}{HTML}{F3F6F6}
\definecolor{warn}{HTML}{A6423A}

\setlength{\parindent}{0pt}
\setlength{\parskip}{6pt}
\renewcommand{\arraystretch}{1.35}

\pagestyle{fancy}
\fancyhf{}
\lhead{\small\textcolor{textgray}{Чек-лист ЕГЭ}}
\rhead{\small\textcolor{textgray}{Тимофей Васильченко}}
\cfoot{\small\thepage}
\renewcommand{\headrulewidth}{0.4pt}
\renewcommand{\headrule}{\hbox to\headwidth{\color{accent}\leaders\hrule height \headrulewidth\hfill}}

\setlist[itemize]{leftmargin=6mm,itemsep=2pt,topsep=2pt}
\setlist[enumerate]{leftmargin=7mm,itemsep=4pt,topsep=4pt}

\newcommand{\sectionline}{\vspace{1mm}{\color{accent}\hrule height 0.6pt}\vspace{2mm}}
\newcommand{\important}[1]{\textcolor{darkaccent}{\textbf{#1}}}
\newcommand{\warning}[1]{\textcolor{warn}{\textbf{#1}}}
\newcommand{\tg}{\operatorname{tg}}
\newcommand{\ctg}{\operatorname{ctg}}
\newcommand{\arctg}{\operatorname{arctg}}
\newcommand{\Z}{\mathbb{Z}}

\begin{document}

\begin{titlepage}
\thispagestyle{empty}
\centering

\vspace*{28mm}

{\Huge\bfseries\textcolor{darkaccent}{Чек-лист}\par}
\vspace{8mm}

{\Large\bfseries Тригонометрические уравнения:\\[1mm]
формулы приведения, двойной угол, понижение степени и замена\par}

\vspace{18mm}

{\large Лёвин Артём Александрович\\[1mm]
эксклюзивно для Тимофея Васильченко\par}

\vspace{10mm}

{\large \today\par}

\vfill

\begin{tikzpicture}[remember picture,overlay]
\draw[accent,line width=1.1pt]
  ([xshift=18mm,yshift=26mm]current page.south west)
  .. controls ([xshift=55mm,yshift=46mm]current page.south west)
  and ([xshift=-70mm,yshift=8mm]current page.south east)
  ..
  ([xshift=-18mm,yshift=30mm]current page.south east);

\draw[darkaccent,line width=0.45pt]
  ([xshift=28mm,yshift=18mm]current page.south west)
  .. controls ([xshift=75mm,yshift=34mm]current page.south west)
  and ([xshift=-65mm,yshift=20mm]current page.south east)
  ..
  ([xshift=-30mm,yshift=15mm]current page.south east);
\end{tikzpicture}

\vspace*{18mm}
\end{titlepage}

\tableofcontents
\newpage

\section{Главный принцип тригонометрического уравнения}
\sectionline

\important{Главный императив.}
Перед тем как решать тригонометрическое уравнение, нужно стремиться к двум целям:
\[
\text{один угол}
\qquad\text{и}\qquad
\text{одна функция}.
\]

То есть выражения вида
\[
\sin 2x,\qquad \cos x,\qquad \sin\left(x+\frac{\pi}{2}\right)
\]
желательно преобразовать так, чтобы в уравнении появились одинаковые углы и одинаковые функции.

\important{Типовая схема.}
\begin{enumerate}
\item Упростить сдвиги углов через формулы приведения или формулы суммы/разности.
\item Раскрыть двойные углы, если встречаются \(2x\), \(3x\), \(\frac{x}{2}\) и другие связанные углы.
\item Добиться повторяющегося выражения.
\item Сделать замену.
\item Решить алгебраическое уравнение.
\item Вернуться к тригонометрии и решить простейшие уравнения.
\end{enumerate}

\warning{Типичная ошибка.}
Нельзя начинать с хаотического применения всех известных формул. Сначала нужно понять, какой угол и какую функцию выгодно оставить.

\section{Формулы суммы и приведения}
\sectionline

\important{Формула синуса суммы.}
\[
\sin(\alpha+\beta)=\sin\alpha\cos\beta+\cos\alpha\sin\beta.
\]

\important{Формула синуса разности.}
\[
\sin(\alpha-\beta)=\sin\alpha\cos\beta-\cos\alpha\sin\beta.
\]

\important{Формула косинуса разности.}
\[
\cos(\alpha-\beta)=\cos\alpha\cos\beta+\sin\alpha\sin\beta.
\]

\important{Пример.}
\[
\sin\left(x+\frac{\pi}{2}\right)
=
\sin x\cos\frac{\pi}{2}+\cos x\sin\frac{\pi}{2}.
\]
Так как
\[
\cos\frac{\pi}{2}=0,\qquad \sin\frac{\pi}{2}=1,
\]
то
\[
\sin\left(x+\frac{\pi}{2}\right)=\cos x.
\]

\important{Ещё один пример.}
\[
\sin\left(\frac{7\pi}{2}-x\right).
\]
Так как
\[
\sin\frac{7\pi}{2}=-1,\qquad \cos\frac{7\pi}{2}=0,
\]
то
\[
\sin\left(\frac{7\pi}{2}-x\right)
=
\sin\frac{7\pi}{2}\cos x-\cos\frac{7\pi}{2}\sin x
=
-\cos x.
\]
Следовательно,
\[
\sin^2\left(\frac{7\pi}{2}-x\right)=\cos^2x.
\]

\warning{Типичная ошибка.}
Если выражение затем возводится в квадрат, минус может исчезнуть:
\[
(-\cos x)^2=\cos^2x.
\]
Но до возведения в квадрат знак терять нельзя.

\section{Формулы двойного угла}
\sectionline

\important{Косинус двойного угла имеет три формы.}
\[
\cos 2\alpha=\cos^2\alpha-\sin^2\alpha,
\]
\[
\cos 2\alpha=2\cos^2\alpha-1,
\]
\[
\cos 2\alpha=1-2\sin^2\alpha.
\]

\important{Как выбрать формулу.}
Выбирайте ту форму, которая делает функции одинаковыми.

\begin{itemize}
\item Если рядом есть \(\cos x\) или \(\cos^2x\), выгодна формула
\[
\cos 2x=2\cos^2x-1.
\]
\item Если рядом есть \(\sin x\) или \(\sin^2x\), выгодна формула
\[
\cos 2x=1-2\sin^2x.
\]
\item Если нужно узнать разность квадратов, полезна форма
\[
\cos 2x=\cos^2x-\sin^2x.
\]
\end{itemize}

\important{Пример.}
Решим уравнение
\[
\cos 2x=\sin\left(x+\frac{\pi}{2}\right).
\]
Сначала:
\[
\sin\left(x+\frac{\pi}{2}\right)=\cos x.
\]
Значит:
\[
\cos 2x=\cos x.
\]
Так как справа стоит \(\cos x\), раскрываем через косинус:
\[
2\cos^2x-1=\cos x.
\]
Переносим всё влево:
\[
2\cos^2x-\cos x-1=0.
\]
Делаем замену:
\[
t=\cos x,\qquad -1\le t\le1.
\]
Получаем:
\[
2t^2-t-1=0.
\]
Корни:
\[
t=1,\qquad t=-\frac12.
\]
Возвращаемся:
\[
\cos x=1
\quad\text{или}\quad
\cos x=-\frac12.
\]

\warning{Типичная ошибка.}
Если в уравнении есть \(\cos 2x\), не надо автоматически раскрывать его первой попавшейся формулой. Нужна та форма, которая приводит к одной функции.

\section{Формулы понижения степени}
\sectionline

\important{Основные формулы.}
\[
\cos^2\alpha=\frac{1+\cos 2\alpha}{2},
\]
\[
\sin^2\alpha=\frac{1-\cos 2\alpha}{2}.
\]

\important{Когда они нужны.}
Формулы понижения степени нужны, когда после преобразований возникает уравнение вида
\[
\sin^2x=\frac14,
\qquad
\cos^2x=\frac12,
\]
а дальше желательно перейти к уравнению с \(\cos 2x\).

\important{Пример.}
Пусть после преобразований получилось:
\[
\sin^2x=\frac14.
\]
Понижаем степень:
\[
\frac{1-\cos 2x}{2}=\frac14.
\]
Умножаем на \(2\):
\[
1-\cos 2x=\frac12.
\]
\[
-\cos 2x=-\frac12.
\]
\[
\cos 2x=\frac12.
\]

\warning{Типичная ошибка.}
При понижении степени угол удваивается:
\[
\sin^2x \longrightarrow \cos 2x.
\]
Нельзя писать \(\cos x\) вместо \(\cos 2x\).

\section{Замена: что именно обозначать за \(t\)}
\sectionline

\important{Главное правило.}
За \(t\) обозначают не обязательно \(\sin x\) или \(\cos x\), а то выражение, которое действительно повторяется.

\important{Пример 1.}
\[
4\cos^4x-4\cos^2x+1=0.
\]
Здесь повторяется не \(\cos x\), а \(\cos^2x\). Поэтому:
\[
t=\cos^2x,\qquad 0\le t\le1.
\]
Тогда:
\[
4t^2-4t+1=0.
\]
Сворачиваем:
\[
(2t-1)^2=0.
\]
Значит,
\[
t=\frac12.
\]
Возвращаемся:
\[
\cos^2x=\frac12.
\]

\important{Пример 2.}
\[
4\sin^4x-4\sin^2x+1=0.
\]
Здесь:
\[
t=\sin^2x,\qquad 0\le t\le1.
\]
Получаем:
\[
4t^2-4t+1=0.
\]

\warning{Типичная ошибка.}
Если заменить \(\cos x=t\) в уравнении с \(\cos^4x\) и \(\cos^2x\), получится уравнение четвёртой степени. Гораздо проще заменить \(\cos^2x=t\).

\section{ОДЗ и область значений при замене}
\sectionline

\important{При замене обязательно думайте об ограничениях.}

Если
\[
t=\sin x
\quad\text{или}\quad
t=\cos x,
\]
то
\[
-1\le t\le1.
\]

Если
\[
t=\sin^2x
\quad\text{или}\quad
t=\cos^2x,
\]
то
\[
0\le t\le1.
\]

\important{Зачем это нужно.}
После замены могут появиться алгебраические корни, которые не имеют смысла для тригонометрии.

\important{Пример.}
Пусть получилось:
\[
t^2-3t+2=0,
\qquad
t=\sin^2x.
\]
Корни:
\[
t=1,\qquad t=2.
\]
Но
\[
0\le \sin^2x\le1.
\]
Следовательно,
\[
t=2
\]
не подходит.

\warning{Типичная ошибка.}
Если корень не подходит по области значений \(t\), нужно явно написать:
\[
t=2\quad \text{не подходит, так как }0\le t\le1.
\]

\section{Дробные выражения: замена обратной функции}
\sectionline

Иногда в уравнении появляются дроби:
\[
\frac{4}{\cos^2x}-\frac{11}{\cos x}+6=0.
\]

\important{Выгодная замена.}
В такой ситуации удобно обозначить:
\[
t=\frac1{\cos x}.
\]
Тогда:
\[
\frac{1}{\cos^2x}=t^2.
\]
Уравнение превращается в квадратное:
\[
4t^2-11t+6=0.
\]

\important{ОДЗ.}
Так как есть деление на \(\cos x\), обязательно:
\[
\cos x\ne0.
\]

\important{Корни квадратного уравнения.}
\[
4t^2-11t+6=0.
\]
\[
D=121-96=25.
\]
\[
t_1=\frac{11-5}{8}=\frac34,
\qquad
t_2=\frac{11+5}{8}=2.
\]
Но
\[
t=\frac1{\cos x}.
\]
Значит:
\[
\frac1{\cos x}=\frac34
\quad\text{или}\quad
\frac1{\cos x}=2.
\]
Первое уравнение даёт
\[
\cos x=\frac43,
\]
что невозможно. Второе даёт
\[
\cos x=\frac12.
\]

\warning{Типичная ошибка.}
Если замена сложная, область значений \(t\) можно не выписывать заранее, но после обратной замены нужно обязательно проверить невозможные случаи.

\section{Узнавание формул в непривычном виде}
\sectionline

\important{Формулы часто спрятаны.}
Например, выражение
\[
\cos\frac{7x}{2}\cos\frac{x}{2}+\sin\frac{7x}{2}\sin\frac{x}{2}
\]
нужно узнать как косинус разности:
\[
\cos\frac{7x}{2}\cos\frac{x}{2}+\sin\frac{7x}{2}\sin\frac{x}{2}
=
\cos\left(\frac{7x}{2}-\frac{x}{2}\right)
=
\cos 3x.
\]

\important{Пример.}
Рассмотрим уравнение
\[
\cos\frac{7x}{2}\cos\frac{x}{2}+\sin\frac{7x}{2}\sin\frac{x}{2}
=
\cos^2 3x.
\]
Левая часть:
\[
\cos\left(\frac{7x}{2}-\frac{x}{2}\right)=\cos3x.
\]
Получаем:
\[
\cos3x=\cos^2 3x.
\]
Переносим:
\[
\cos^2 3x-\cos3x=0.
\]
Выносим общий множитель:
\[
\cos3x(\cos3x-1)=0.
\]
Значит:
\[
\cos3x=0
\quad\text{или}\quad
\cos3x=1.
\]

\warning{Типичная ошибка.}
Если формула дана в справочных материалах, это не значит, что её легко увидеть. Нужно тренироваться узнавать структуру:
\[
\cos A\cos B+\sin A\sin B=\cos(A-B).
\]

\section{Когда не стоит сразу извлекать корень}
\sectionline

\important{Опасный путь.}
Если получилось
\[
\sin^2x=\frac14,
\]
можно написать:
\[
\sin x=\pm\frac12.
\]
Это верно, но дальше появятся две серии решений.

\important{Почему иногда лучше понижать степень.}
Если впереди пункт отбора корней на промежутке, то две серии могут усложнить работу. В таких задачах часто выгоднее:
\[
\sin^2x=\frac14
\]
заменить на:
\[
\frac{1-\cos2x}{2}=\frac14.
\]
Тогда получится одно уравнение:
\[
\cos2x=\frac12.
\]

\warning{Типичная ошибка.}
Извлечение корня не является ошибкой, но может сделать пункт отбора корней намного тяжелее. Выбирайте путь с учётом следующего пункта задачи.

\section{Короткий итоговый чек-лист}
\sectionline

\begin{longtable}{p{0.31\linewidth}p{0.63\linewidth}}
\toprule
\textbf{Ситуация} & \textbf{Что делать} \\
\midrule
Разные углы & Приводить к одному углу через формулы приведения, суммы, разности или двойного угла. \\
Разные функции & Выбирать такую формулу, чтобы осталась одна функция: \(\sin\) или \(\cos\). \\
Есть \(\cos2x\) & Раскрывать как \(2\cos^2x-1\), \(1-2\sin^2x\) или \(\cos^2x-\sin^2x\) по ситуации. \\
Есть квадрат функции & Можно использовать формулы понижения степени. \\
Есть \(\sin^4x,\cos^4x\) & Часто замена: \(t=\sin^2x\) или \(t=\cos^2x\). \\
Деление на \(\cos x\) или \(\sin x\) & Не забывать ОДЗ: знаменатель не равен нулю. \\
Есть \(\frac1{\cos x}\) & Возможна замена \(t=\frac1{\cos x}\). \\
После замены & Обязательно указать область значений \(t\), если она простая. \\
Корень не подходит & Явно написать, почему он не подходит. \\
Пункт с отбором корней & Иногда лучше понижать степень, а не сразу писать \(\pm\). \\
Формула спрятана & Узнавать структуры вида \(\cos A\cos B+\sin A\sin B\). \\
\bottomrule
\end{longtable}

\newpage

\section{Тренировочный блок}
\sectionline

\textbf{Инструкция.} Решите задачи самостоятельно. Ответы находятся на следующей странице.

\subsection*{Средний уровень}

\begin{enumerate}
\item Приведите уравнение к квадратному относительно \(t=\cos x\):
\[
\cos2x=\sin\left(x+\frac{\pi}{2}\right).
\]

\item Приведите уравнение к уравнению вида \(\cos2x=c\):
\[
\sin^2x=\frac14.
\]

\item Решите уравнение:
\[
4\cos^4x-4\cos^2x+1=0.
\]
\end{enumerate}

\subsection*{Выше среднего уровня}

\begin{enumerate}[resume]
\item Решите уравнение:
\[
4\sin^4x-4\sin^2x+1=0.
\]

\item Решите уравнение:
\[
\frac{4}{\cos^2x}-\frac{11}{\cos x}+6=0.
\]

\item Упростите левую часть и решите уравнение:
\[
\cos\frac{7x}{2}\cos\frac{x}{2}+\sin\frac{7x}{2}\sin\frac{x}{2}
=
\cos^2 3x.
\]

\item Решите уравнение:
\[
\cos2x+\sin^2x=\frac12.
\]

\item Решите уравнение:
\[
\sin\left(\frac{5\pi}{2}-x\right)=\cos2x.
\]

\item Решите уравнение:
\[
\sin^2\left(\frac{7\pi}{2}-x\right)-\cos2x=0.
\]
\end{enumerate}

\subsection*{Сложный уровень}

\begin{enumerate}[resume]
\item Доведите уравнение до квадратного относительно одной тригонометрической функции и решите его:
\[
\sin^2\frac{x}{2}-\cos^2\frac{x}{2}=\cos x.
\]
\end{enumerate}

\newpage

\section{Ответы}
\sectionline

\begin{table}[h]
\centering
\caption{Ответы к тренировочному блоку}
\begin{tabularx}{\linewidth}{cX}
\toprule
\textbf{№} & \textbf{Ответ} \\
\midrule
1 & \(2t^2-t-1=0,\quad t=\cos x\). \\
2 & \(\cos2x=\dfrac12\). \\
3 & \(x=\pm\dfrac{\pi}{4}+\pi k,\quad k\in\Z\). \\
4 & \(x=\dfrac{\pi}{4}+\dfrac{\pi k}{2},\quad k\in\Z\). \\
5 & \(x=\pm\dfrac{\pi}{3}+2\pi k,\quad k\in\Z\). \\
6 & \(\cos3x=0\) или \(\cos3x=1\); то есть \(x=\dfrac{\pi}{6}+\dfrac{\pi k}{3}\) или \(x=\dfrac{2\pi k}{3}\), \(k\in\Z\). \\
7 & \(x=\pm\dfrac{\pi}{3}+2\pi k,\quad k\in\Z\). \\
8 & \(x=2\pi k\) или \(x=\pm\dfrac{2\pi}{3}+2\pi k,\quad k\in\Z\). \\
9 & \(x=\dfrac{\pi}{4}+\dfrac{\pi k}{2},\quad k\in\Z\). \\
10 & Все \(x\in\mathbb R\), для которых выражения определены; после преобразования получается тождество \(-\cos x=\cos x\), значит \(\cos x=0\), то есть \(x=\dfrac{\pi}{2}+\pi k,\quad k\in\Z\). \\
\bottomrule
\end{tabularx}
\end{table}

\end{document}